% !TeX root = ../初等数学笔记.tex
% 由主文档通过 \input 载入；请编译主文档。
\setcounter{chapter}{16} % 保留原书章号。
\chapter{角与三角函数}
\label{chap:em-trigonometry}
三角函数把角度与长度、坐标联系起来。单位圆定义统一了锐角、负角与超过一周的角，也使周期、奇偶性和基本恒等式成为同一几何结构的不同表现。

\section{有向角、弧度与单位圆}
以逆时针旋转为正，顺时针为负，同一终边可对应相差整周的角。\emph{弧度}（radian）把角度定义为所截有向弧长与半径之比，一周为 $2\pi$，故 $180^\circ=\pi$ 弧度。除特别标注角度符号外，本书三角表达式中的角均用弧度。

半径为 $r>0$、非负圆心角 $0\le\theta\le2\pi$ 的普通扇形，其弧长与面积为
\[
 s=r\theta,\qquad S=\frac12r^2\theta.
\]
这里 $s$ 是非负几何长度；对带符号的旋转角，方向与弧长大小应分别处理。多次绕行所积累的弧长也不等于终边间较短圆弧的长度。

\section{三角函数的单位圆定义}
\begin{definition}
从 $(1,0)$ 绕原点旋转角 $x$ 到单位圆上的点 $P$，其坐标写为 $(\cos x,\sin x)$，分别定义\emph{余弦}（cosine）和\emph{正弦}（sine）。当 $\cos x\ne0$ 时，定义\emph{正切}（tangent）$\tan x=\sin x/\cos x$。
\end{definition}
同理 $\cot x=\cos x/\sin x$、$\sec x=1/\cos x$、$\csc x=1/\sin x$，均要求分母非零。正弦与余弦在整个实轴定义；正切排除 $x=\pi/2+k\pi$，余切排除 $x=k\pi$，其中 $k\in\mathbb Z$。

对锐角，把相应直角三角形按斜边缩放到 $1$，即可看到 $\sin x$ 为对边与斜边之比、$\cos x$ 为邻边与斜边之比、$\tan x$ 为对边与邻边之比。相似性保证这些比值与三角形大小无关。

\section{定义域、符号与基本恒等式}
单位圆方程给出
\begin{equation}\label{eq:em-trig-pythagorean}
 \sin^2x+\cos^2x=1.
\end{equation}
分别除以允许非零的 $\cos^2x$ 或 $\sin^2x$，得到 $1+\tan^2x=\sec^2x$、$1+\cot^2x=\csc^2x$。记号 $\sin^2x$ 表示 $(\sin x)^2$，不是 $\sin(x^2)$。

四个象限的正弦符号依次为 $+,+,-,-$，余弦为 $+,-,-,+$，正切由二者之商确定。坐标轴上的值直接按坐标计算。等边三角形作高与等腰直角三角形分别给出
\[
 \sin\frac\pi6=\frac12,\quad\cos\frac\pi6=\frac{\sqrt3}2,\quad
 \sin\frac\pi4=\cos\frac\pi4=\frac{\sqrt2}2,\quad
 \sin\frac\pi3=\frac{\sqrt3}2,\quad\cos\frac\pi3=\frac12.
\]
若由 $\sin x$ 求 $\cos x$，平方关系只能给出 $\cos x=\pm\sqrt{1-\sin^2x}$，符号必须由角的范围确定。

\section{诱导关系、周期与图像}
单位圆关于横轴、原点及直线 $y=x$ 的对称给出
\[
 \sin(-x)=-\sin x,\quad \cos(-x)=\cos x,\quad
 \sin(x+\pi)=-\sin x,\quad\cos(x+\pi)=-\cos x,
\]
以及 $\sin(\pi/2-x)=\cos x$。转过整周后位置不变，所以正弦、余弦以 $2\pi$ 为周期，正切以 $\pi$ 为周期；它们分别没有更小的正周期，可由零点位置与相邻区间的符号验证。

在半周 $[-\pi/2,\pi/2]$ 上，单位圆纵坐标随角严格增加，故正弦严格递增；余弦在 $[0,\pi]$ 上严格递减。结合对称与周期便能确定全部单调区间。图像中正弦与余弦的值域为 $[-1,1]$，正切在每个 $(k\pi-\pi/2,k\pi+\pi/2)$ 上严格递增并覆盖实轴。一般周期函数可能没有最小正周期，例如常函数的每个正数都是周期。

\section{图像参数与反三角函数选读}
对 $y=A\sin(\omega x+\varphi)+d$，若 $A\ne0,\omega\ne0$，振幅为 $|A|$，最小正周期为 $2\pi/|\omega|$，中线为 $y=d$。相位应通过 $\omega(x+\varphi/\omega)$ 解释水平位移，不能把 $\varphi$ 本身直接当作平移长度。若 $A=0$ 或 $\omega=0$，函数退化为常量，前述最小周期公式不适用。

把正弦限制在 $[-\pi/2,\pi/2]$，余弦限制在 $[0,\pi]$，正切限制在 $(-\pi/2,\pi/2)$，它们均单射，分别定义反正弦 $\arcsin$、反余弦 $\arccos$ 和反正切 $\arctan$。反正弦、反余弦输入在 $[-1,1]$，反正切输入可为任意实数。主值只是指定区间内的一个角，不能代表三角方程的全部解。

\section{例题、习题与学习检查}
\begin{example}
已知 $x$ 的终边在第二象限，且 $\sin x=3/5$，求 $\cos x,\tan x$。
\end{example}
\begin{solve}
由式 \eqref{eq:em-trig-pythagorean} 有 $\cos^2x=16/25$，第二象限余弦为负，所以 $\cos x=-4/5$，进而 $\tan x=-3/4$。象限条件决定了平方根的符号。
\end{solve}
\begin{enumerate}
\item 将 $150^\circ$ 化为弧度并求其正弦、余弦。答：$5\pi/6,1/2,-\sqrt3/2$。
\item 求 $2\sin(3x-\pi)+1$ 的振幅、周期和值域。答：$2,2\pi/3,[-1,3]$。
\item 求 $\arcsin(\sin(5\pi/6))$。答：$\pi/6$，反函数只反转所选单射分支。
\end{enumerate}

